\documentclass[11pt]{article}

\usepackage[margin=1in]{geometry}
\usepackage{amsmath}
\usepackage{amssymb}
\usepackage{graphicx}
\usepackage{float}
\usepackage{booktabs}
\usepackage{xcolor}
\usepackage{hyperref}
\usepackage{enumitem}
\usepackage{caption}
\usepackage{listings}
\usepackage{tikz}
\usetikzlibrary{arrows.meta, positioning, fit, backgrounds, calc, shapes.geometric}

\hypersetup{colorlinks=true, linkcolor=blue, urlcolor=blue, citecolor=blue}

\lstset{
  basicstyle=\ttfamily\footnotesize,
  breaklines=true,
  frame=single,
  rulecolor=\color{gray!30},
  backgroundcolor=\color{gray!5},
  numbers=none,
}

\title{\textbf{Closed-Loop Reactive Control for On-the-Move Grasping}\\[4pt]
\large Progress writeup --- toolchain validation and reference QP on Apple Silicon}
\author{Anurag Paresh Shetye \quad\textbar\quad ARMS Lab, IIT Bombay (visiting) \quad\textbar\quad collab.\ with University of Osaka}
\date{31 May 2026}

\begin{document}
\maketitle

\section{Goal}

The system extends the open-loop pipeline of Kiyokawa et al.\
(IEEE/SICE SII 2025) for on-the-move grasping by introducing a high-rate
reactive controller that consumes the predicted grasp pose
$(x, y, \theta)$ in world frame and continuously recomputes arm joint
velocities against a fused live base pose estimate, closing the loop on
velocity-tracking error during the approach. The platform is the
TurtleBot 3 Waffle (differential drive) with an OpenManipulator-X
(4-DOF arm, $380\,$mm reach), a strictly more arm-constrained setup than
the Toyota HSR used in the original paper.

The architectural target is the figure shared on 20 May 2026; the
controller follows the Haviland-Corke Holistic formulation
(IEEE RA-L 2022), adapted to the constrained 4-DOF arm and integrated
with a swappable grasp predictor (FCN reimplementation, GG-CNN, or
OpenVLA-OFT). This document records the first concrete implementation
step, taken on a development laptop ahead of moving to the lab Linux
machine and physical hardware.

\section{Setup}

All work in this writeup runs on a MacBook Pro M3 Pro (18\,GB RAM,
Apple Silicon), with Python 3.9 in a project-local virtual environment.
The intent of working on macOS first is to verify the controller-side
toolchain in isolation, without any simulator or ROS runtime
dependencies, so that subsequent integration on the lab Linux machine
addresses one new concern at a time.

\paragraph{Repository.}
\begin{itemize}[noitemsep]
  \item \texttt{github.com/orionop/mobile-grasping}
  \item Layout: \texttt{src/} (controller, predictor, pipeline modules),
        \texttt{Docs/} (this writeup, progress log),
        \texttt{tests/} (unit tests for the toolchain),
        \texttt{experiments/} (numbered experiment scripts),
        \texttt{notebooks/} (math derivations and sanity checks).
  \item Dependencies pinned in \texttt{pyproject.toml}; numpy held at
        \texttt{<2} because the \texttt{roboticstoolbox-python} C
        extension is built against numpy~1.x.
\end{itemize}

\paragraph{Toolchain.} The implementation builds on three open-source
components:
\begin{itemize}[noitemsep]
  \item \texttt{roboticstoolbox-python}~1.1.1 (Corke, Haviland) ---
        kinematic models, Jacobians, manipulator URDFs.
  \item \texttt{spatialmath-python}~1.1 --- \texttt{SE(3)} algebra used
        for end-effector pose representation.
  \item \texttt{qpsolvers}~4.8 with \texttt{osqp}~1.1 backend --- the QP
        solver used to compute joint velocities each control step.
\end{itemize}

\section{Code}

The first experiment, \texttt{experiments/exp01\_holistic\_reference.py},
implements one control step of the Holistic QP on the reference Frankie
platform (Franka Panda 7-DOF arm). The intent is purely to validate the
solver pipeline end-to-end on Apple Silicon before adapting any of the
formulation to the constrained 4-DOF arm.

\paragraph{Decision variable.} Joint velocities for the mobile
manipulator, augmented with an end-effector slack vector:
\[
  x = \begin{bmatrix} \dot{q}_{\text{base}} \\
                       \dot{q}_{\text{arm}} \\
                       \delta \end{bmatrix}
     \in \mathbb{R}^{2 + 7 + 6} = \mathbb{R}^{15}.
\]

\paragraph{Quadratic program.}
\[
  \min_{x}\ \tfrac{1}{2}\, x^\top Q\, x + C^\top x
  \quad\text{s.t.}\quad
  J\, x = \nu^*_e,\quad
  X^{-} \le x \le X^{+},
\]
where the augmented Jacobian is
\(
  J = \begin{bmatrix} J_{\text{base}} & J_{\text{arm}}(q) & \mathbb{1}_{6 \times 6}\end{bmatrix}
\)
and the slack penalty in $Q$ is set high (1\,e\,6) so that the equality
constraint is enforced tightly in this reference test; the full
adaptive-weight scheme from RA-L 2022 is left for the next iteration.
The arm joint velocities are bounded to $\pm 1.5\,$rad/s and the base
virtual joints to $\pm 0.3\,$m/s.

\paragraph{Test case.} The Panda is placed in the toolbox-provided
\texttt{qr} ready configuration. The desired end-effector twist in the
world frame is
\(\nu^*_e = (0.1,\ 0,\ 0,\ 0,\ 0,\ 0)\,\)m/s, i.e.\ pure $+x$
translation with no rotation, no $y/z$ motion. The QP is solved once and
the realised twist is reconstructed as
\(\nu^{\text{realised}}_e = J_{\text{arm}}(q)\,\dot{q}_{\text{arm}}\)
to check tracking quality.

\section{Result}

\begin{table}[H]
\centering
\begin{tabular}{lr}
\toprule
\textbf{Quantity} & \textbf{Value} \\
\midrule
End-effector start position (Panda \texttt{qr}) & $(0.484,\ 0.000,\ 0.413)\,$m \\
Desired end-effector twist                       & $(0.1,\ 0,\ 0,\ 0,\ 0,\ 0)\,$m/s \\
Base joint velocities returned                   & $(0,\ 0)$ \\
Arm joint velocities returned (rad/s)            & $(0,\ 0.308,\ 0,\ 0.297,\ 0,\ 0.012,\ 0)$ \\
Slack returned (6-D)                              & $\approx 0$ \\
Realised end-effector twist                      & $(0.1,\ 0,\ 0,\ 0,\ 0,\ 0)\,$m/s \\
Tracking error (L2 norm)                          & $2\times 10^{-6}$ \\
\bottomrule
\end{tabular}
\caption{One-step output of the reference Holistic QP on the Panda.
The solver returns joint velocities that realise the requested
end-effector twist to within $2\times 10^{-6}$, with zero slack absorbed
at the chosen penalty.}
\end{table}

Three observations are worth recording:

\begin{itemize}[noitemsep]
  \item The QP returns a feasible, well-conditioned solution in well
        under $10\,$ms on the M3 Pro, comfortably below the control
        period required for the $\sim 200\,$Hz target rate.
  \item Only three of the seven arm joints (joints 2, 4, 6) carry
        motion. At the \texttt{qr} configuration with a pure $+x$
        translational request, this is the kinematically minimal
        solution and indicates the Jacobian and the cost weighting are
        producing sensible joint allocations.
  \item Zero slack at the chosen 1\,e\,6 penalty confirms that the
        equality constraint is being enforced as intended. When the
        penalty is lowered to 1\,e\,2, the solver trades joint motion
        against tracking error in the expected direction (slack
        absorbs $\sim 1.6\%$ of the desired velocity), which validates
        the Holistic slack mechanism is wired correctly.
\end{itemize}

\section{What this validates and what remains}

The Day 1 result is narrow but iron-clad: the controller-side toolchain
(\texttt{roboticstoolbox-python}, \texttt{qpsolvers}, \texttt{osqp})
works correctly on Apple Silicon, the augmented Jacobian formulation
from Haviland-Corke RA-L 2022 implements cleanly in Python, and the QP
solver produces sensible outputs within a control-rate-compatible
time budget.

\paragraph{Pending on the controller side.} The reference experiment
implements the QP structure only; the next iterations will add the
remaining elements of the Holistic formulation:
\begin{itemize}[noitemsep]
  \item Manipulability gradient in the linear cost vector $C$, using
        the arm-only Yoshikawa measure
        \(\hat{J}_m = [\mathbf{0}_b,\ J^a_m(q_a)]\).
  \item Adaptive cost weights $\lambda_q$ and $\lambda_\delta$ as
        functions of the end-effector pose error magnitude
        $\|e\|$, so that the base contributes more when the goal is
        far and the arm contributes more when it is close.
  \item Velocity-damper inequality constraints for joint position
        limits, with the influence distance $\rho_i$ and minimum
        distance $\rho_s$ as parameters.
  \item Base orientation cost term $\theta_\epsilon$ to keep the base
        facing the end-effector.
  \item Position-Based Servoing wrapper that converts a desired
        end-effector pose into the desired twist $\nu^*_e$ each step.
\end{itemize}

\paragraph{Pending on the platform side.} The reference test uses the
Panda model that ships with the toolbox. The constrained-platform
adaptation will require a 4-DOF arm Jacobian for the OpenManipulator-X
and a differential-drive virtual-joint representation for the TB3
Waffle base. With four arm degrees of freedom there is no redundancy
slack against the full 6-D end-effector twist; the equality constraint
will need to be projected onto the four task-space dimensions that
Kiyokawa's grasp pose actually specifies ($x$, $y$, $z$ implicit, and
$\theta$ around the vertical axis).

\paragraph{Pending on the system side.} The simulator and ROS bring-up
have not yet been started; that work moves to the lab Linux machine in
the coming week. The intended sequence is: load the combined
TB3 + OpenManipulator-X URDF in a simulator with the chosen ROS
distribution; verify base-only and arm-only motion with stock
controllers; integrate the QP controller as a ROS node that consumes a
grasp pose and a live base pose estimate and publishes arm joint
velocity commands; run a first open-loop versus closed-loop comparison
on a single object at low base velocity (0.10\,m/s) to establish a
baseline before sweeping velocities.

\section{Repository status}

\begin{itemize}[noitemsep]
  \item \texttt{src/} --- intentionally empty so far; the modular QP
        solver will land here in the next commit, with the experiment
        script reduced to a thin wrapper around the module so the same
        solver code is reused by the ROS node.
  \item \texttt{experiments/exp01\_holistic\_reference.py} --- the
        reference QP described above.
  \item \texttt{tests/test\_imports.py} --- a small sanity test
        verifying that all toolchain components import and operate on
        Apple Silicon.
  \item \texttt{Docs/writeup.tex}, \texttt{Docs/writeup.pdf} --- this
        document.
  \item \texttt{Docs/progress\_log.md} --- internal running log of
        milestones, complementary to this writeup.
  \item \texttt{pyproject.toml}, \texttt{.gitignore} --- standard.
\end{itemize}

\end{document}
